\begin{figure}[htbp]
    \centering
    % Ensure you have \usepackage{tikz} and \usetikzlibrary{positioning, arrows.meta, calc} in your preamble
    \begin{tikzpicture}[
        >=Stealth,
        node distance=2cm and 2.5cm,
        font=\sffamily\small,
        % Define block styles
        block/.style={rectangle, draw, rounded corners, thick, minimum width=2.6cm, minimum height=1.5cm, align=center},
        E_node/.style={block, fill=blue!5, draw=blue!70!black},
        I_node/.style={block, fill=red!5, draw=red!70!black},
        io/.style={circle, draw, thick, fill=gray!5, minimum size=1.0cm, align=center},
        % Arrow styles
        line/.style={draw, thick, ->},
        inh_line/.style={draw, thick, ->, red!80!black}
    ]

        % --- Nodes ---
        \node[io] (in) {$x[n]$};
        
        \node[E_node, right=of in] (E) {
            \textbf{Excitatory (E)} \\ 
            Population State \\ 
            $E[n]$
        };
        
        \node[I_node, below=1.8cm of E] (I) {
            \textbf{Inhibitory (I)} \\ 
            Quadrature State \\ 
            $I[n]$
        };
        
        \node[io, right=of E] (out) {$O[n]$};

        % --- Main Feedforward Path ---
        \path[line] (in) -- node[above] {$g$} (E);
        \path[line] (E) -- node[above] {spikes} (out);
        
        % --- Coupling Loops (The Rotation Matrix) ---
        % Forward Excitation: E -> I
        \path[line] ($(E.south east)!0.3!(E.south)$) -- 
            node[right, pos=0.5] {$+p\sin(\phi)$}
            ($(I.north east)!0.3!(I.north)$);
            
        % Delayed Negative Feedback: I -> E
        \path[inh_line] ($(I.north west)!0.3!(I.north)$) -- 
            node[left, pos=0.5] {$-p\sin(\phi)$} 
            ($(E.south west)!0.3!(E.south)$);
            
        % --- Self-Decay Loops (Damping terms) ---
        % E Self-loop
        \draw[line] ([xshift=+0.25cm]E.north) .. controls +(0.4, 0.8) and +(-0.4, 0.8) .. ([xshift=-0.25cm]E.north)
            node[above, midway] {$p\cos(\phi)$};
            
        % I Self-loop
        \draw[line] ([xshift=+0.25cm]I.south) .. controls +(0.4, -0.8) and +(-0.4, -0.8) .. ([xshift=-0.25cm]I.south)
            node[below, midway] {$p\cos(\phi)$};

    \end{tikzpicture}
    \caption{Equivalent excitatory/inhibitory population interpretation of the implemented cochlear resonator. The recurrent E/I loop is a state-space interpretation of the second-order IIR resonator. Self-decay is scaled by $p\cos(\phi)$, while balanced cross-coupling scales with $\pm p\sin(\phi)$. Together these terms reproduce the same complex-conjugate pole pair as the direct-form filter.}
    \label{fig:ei_population_resonator}
\end{figure}
